% ============================================================
%  附录F 相干态
%  底本：高等教育出版社《固体理论》第二版；原书印刷页 554–556
%  （PDF 页 = 印刷页 + 20）
%  公式编号自检：(F.1)–(F.16)，自动编号 16 个，无缺号。
%
%  已修正的 OCR/原书错误（除全角句点转"."、删 ~ 间隔符、拆开
%  行内公式与中文等常规项外）：
%  1. 原书 p554 (F.8) 第一行印 e^{-αâ⁺}|0>，与其第二行级数
%     Σαⁿ/√n!|n> 及本征方程(F.7)矛盾（e^{-αâ⁺}|0> 的展开为
%     Σ(−α)ⁿ/√n!|n>，本征值应为 −α），系原书排印漏正号，已按
%     物理改排为 e^{αâ⁺}|0> 并在此存照。
%  2. p555 段首"相于态"→"相干态"（OCR 误识）。
%  3. 原书 p555 "它代最小的不确定态"漏"表"字，按文义补为
%     "它代表最小的不确定态"（原书排印漏字，存照）。
%  4. "现在对算符 x̂ 和 p̂ 作相干态平均"OCR 语序错乱按 p555
%     复原；(F.9) 前文"非厄米算符 â"等处 OCR 重复/错位已清理。
%  5. (F.16) 后一段"非厄米算符 â 的非线性方程表示为"OCR 把 â
%     衍入"方程"一词中间，按 p555–556 复原；跨页断词"相干态方
%     法"已接排。
%  6. (F.1) 原书印 a⁺a（算符无帽），与 (F.2)(F.3) 的 â⁺â 不一致，
%     系原书照排，未改。
%  章末按原书照录参考书一句（倪光炯、陈苏卿《高等量子力学》）。
%  遗留问题：无
% ============================================================
\chapter{相干态}

频率为 $\omega$ 的振子哈密顿量为
\begin{equation}
  H_{0}=\hbar\omega\left(a^{+}a+\frac{1}{2}\right)
\end{equation}
振动量子的数算符定义为
\begin{equation}
  \hat{n}\equiv\hat{a}^{+}\hat{a}
\end{equation}
其本征方程可写成
\begin{equation}
  \hat{a}^{+}\hat{a}|n\rangle=n|n\rangle
  \qquad(n=0,1,2,\cdots)
\end{equation}
$|n\rangle$ 代表有确定量子数 $n$ 的状态，有关系式
\begin{equation}
\begin{gathered}
  |n\rangle=\frac{1}{\sqrt{n!}}
  \left(\hat{a}^{+}\right)^{n}|0\rangle\\
  \hat{a}|n\rangle=\sqrt{n}\,|n-1\rangle
  ,\qquad
  \hat{a}^{+}|n\rangle=\sqrt{n+1}\,|n+1\rangle
\end{gathered}
\end{equation}
态 $|n\rangle$ 所对应的振子波函数
\begin{equation}
  \langle x|n\rangle=\left(\sqrt{\pi}\,2^{n}n!\right)^{-\frac{1}{2}}
  \mathrm{H}_{n}(x)\exp\left(-x^{2}/2\right)
\end{equation}
其中
\begin{equation}
  \mathrm{H}_{n}(x)=(-1)^{n}\mathrm{e}^{x^{2}}
  \left(\frac{\mathrm{d}}{\mathrm{d}x}\right)^{n}
  \mathrm{e}^{-x^{2}}
\end{equation}
是 $n$ 阶厄米(Hermite)多项式.

相干态定义为非厄米算符 $\hat{a}$ 的本征态，它满足方程
\begin{equation}
  \hat{a}|\alpha\rangle=\alpha|\alpha\rangle
\end{equation}
其中 $\alpha$ 为复量.满足上式的相干态 $|\alpha\rangle$ 取下列形式：
\begin{equation}
\begin{aligned}
  |\alpha\rangle
  &=\mathrm{e}^{-|\alpha|^{2}/2}\,
  \mathrm{e}^{\alpha\hat{a}^{+}}|0\rangle\\
  &=\mathrm{e}^{-|\alpha|^{2}/2}\sum_{n=0}^{\infty}
  \frac{\alpha^{n}}{\sqrt{n!}}\,|n\rangle
\end{aligned}
\end{equation}
不难证实，若用算符 $\hat{a}$ 作用于 $|\alpha\rangle$ 将得到本征方程(F.7)
\begin{equation*}
\begin{aligned}
  \hat{a}|\alpha\rangle
  &=\mathrm{e}^{-|\alpha|^{2}/2}\sum_{n=0}^{\infty}
  \frac{\alpha^{n}}{\sqrt{n!}}\,\hat{a}|n\rangle\\
  &=\alpha\,\mathrm{e}^{-|\alpha|^{2}/2}\sum_{n=1}^{\infty}
  \frac{\alpha^{n-1}}{\sqrt{(n-1)!}}\,|n-1\rangle\\
  &=\alpha|\alpha\rangle
\end{aligned}
\end{equation*}
因此，$|\alpha\rangle$ 的表示式(F.8)是非厄米算符 $\hat{a}$ 的本征态，其本征值
\begin{equation}
  \alpha=\langle\alpha|\hat{a}|\alpha\rangle
\end{equation}
应当注意，相干态 $|\alpha\rangle$ 不再是具有确定量子数的状态 $|n\rangle$，而是各种 $n$ 的特殊组合态.式(F.8)的第二个等式表明了在谐振子情况下 $|\alpha\rangle$ 与 $|n\rangle$ 的关系.

相干态的另一个重要特征是，它代表最小的不确定态，也可以说是最经典的量子态.众所周知，任意可观察量 $\hat{A}$ 与 $\hat{B}$ 对于一般态 $|\psi\rangle$ 的量子力学平均值 $\langle\hat{A}\rangle=\langle\psi|\hat{A}|\psi\rangle$，若定义 $\Delta\hat{A}=\hat{A}-\langle\hat{A}\rangle$，则均方偏差 $\langle(\Delta\hat{A})^{2}\rangle=\langle\hat{A}^{2}\rangle-(\langle\hat{A}\rangle)^{2}$，在量子力学中已知，可观察量 $\hat{A}$ 与 $\hat{B}$ 之间存在下列不确定关系：
\begin{equation}
  \langle(\Delta\hat{A})^{2}\rangle
  \langle(\Delta\hat{B})^{2}\rangle
  \geqslant\frac{1}{4}
  \left|\langle[\hat{A},\hat{B}]\rangle\right|^{2}
\end{equation}
对于简谐振子系统 $\hat{H}=\frac{1}{2m}\hat{p}^{2}
+\frac{1}{2}m\omega^{2}\hat{x}^{2}$，其中
\begin{equation}
  \hat{x}=\left(\hbar/2m\omega\right)^{1/2}
  \left(\hat{a}^{+}+\hat{a}\right)
  ,\qquad
  \hat{p}=\mathrm{i}\left(m\hbar\omega/2\right)^{1/2}
  \left(\hat{a}^{+}-\hat{a}\right)
\end{equation}
若取 $\hat{A}=\hat{x}$ 与 $\hat{B}=\hat{p}$，因为 $[\hat{x},\hat{p}]=\mathrm{i}\hbar$，代入式(F.10)可求得简谐振子的不确定关系为
\begin{equation}
  \langle(\Delta\hat{x})^{2}\rangle
  \langle(\Delta\hat{p})^{2}\rangle
  \geqslant\frac{1}{4}\hbar^{2}
\end{equation}
当上式中等号成立时称为最小不确定态.现在对算符 $\hat{x}$ 和 $\hat{p}$ 作相干态平均，利用式(F.7)不难算得
\begin{equation}
\begin{aligned}
  &\langle\alpha|\hat{x}|\alpha\rangle
  =\left(\hbar/2m\omega\right)^{1/2}
  \left(\alpha^{*}+\alpha\right)
  ,\\
  &\langle\alpha|\hat{p}|\alpha\rangle
  =-\mathrm{i}\left(m\hbar\omega/2\right)^{1/2}
  \left(\alpha^{*}-\alpha\right)
\end{aligned}
\end{equation}
相应的
\begin{equation}
\begin{aligned}
  &\langle\alpha|\hat{x}^{2}|\alpha\rangle
  =\left(\hbar/2m\omega\right)
  \left[\left(\alpha^{*}+\alpha\right)^{2}+1\right]\\
  &\langle\alpha|\hat{p}^{2}|\alpha\rangle
  =-\left(m\hbar\omega/2\right)
  \left[\left(\alpha^{*}-\alpha\right)^{2}-1\right]
\end{aligned}
\end{equation}
这时，在相干态中，简谐振子的均方偏差为
\begin{equation}
\begin{aligned}
  &\langle(\Delta\hat{x})^{2}\rangle_{\mathrm{c}}
  \equiv\langle\alpha|(\Delta x)^{2}|\alpha\rangle\\
  &\qquad=\langle\alpha|\hat{x}^{2}|\alpha\rangle
  -\left(\langle\alpha|\hat{x}|\alpha\rangle\right)^{2}
  =\left(\hbar/2m\omega\right)\\
  &\langle(\Delta\hat{p})^{2}\rangle_{\mathrm{c}}
  \equiv\langle\alpha|(\Delta p)^{2}|\alpha\rangle
  =\left(m\hbar\omega/2\right)
\end{aligned}
\end{equation}
最后，求得简谐振子在相干态中的均方偏差满足最小不确定关系：
\begin{equation}
  \langle(\Delta\hat{x})^{2}\rangle_{\mathrm{c}}
  \langle(\Delta\hat{p})^{2}\rangle_{\mathrm{c}}
  =\frac{1}{4}\hbar^{2}
\end{equation}
因此，相干态是最小的不确定态.在非线性问题中，相干态将作为一种表象被应用，这时可利用式(F.9)将非厄米算符 $\hat{a}$ 的非线性方程表示为它的相干态平均值方程.在本书第三章中将应用相干态方法讨论一维铁磁链中的孤子运动问题，这是一种良好的准经典近似方法.

有关相干态的详细讨论，读者可参阅倪光炯、陈苏卿的《高等量子力学》（复旦大学出版社，2000年3月）.
